% Created 2026-11-10 mar 20:06
% Intended LaTeX compiler: pdflatex
\documentclass[a4paper]{article}

\usepackage[utf8]{inputenc}
\usepackage[T1]{fontenc}
\usepackage{graphicx}
\usepackage{longtable}
\usepackage{wrapfig}
\usepackage{rotating}
\usepackage[normalem]{ulem}
\usepackage{amsmath}
\usepackage{amssymb}
\usepackage{capt-of}
\usepackage{hyperref}
\usepackage[spanish]{babel}
\usepackage[usenames,dvipsnames]{color} % Required for custom colors
\renewcommand{\ttdefault}{pcr} % MONOESPACIO CON NEGRIT
\usepackage{lastpage}
\usepackage{listings}
\usepackage{listingsutf8}
\renewcommand{\lstlistingname}{Listado}
\lstset{frame=single,inputencoding=utf8,basicstyle=\scriptsize\ttfamily,showstringspaces=false,numbers=none}
\definecolor{MyDarkGreen}{rgb}{0.0,0.4,0.0} % This is the color used for comments
\lstset{ breaklines=true, postbreak=\mbox{\textcolor{red}{$\hookrightarrow$}\space}, keywordstyle=\bfseries, keywordstyle=[1]\color{Blue}\bfseries,  keywordstyle=[2]\color{Purple}\bfseries,  keywordstyle=[3]\color{Blue}\underbar,   identifierstyle=,   commentstyle=\usefont{T1}{pcr}{m}{sl}\color{MyDarkGreen}\small,   stringstyle=\color{Purple},   showstringspaces=false,   tabsize=2,   morecomment=[l][\color{Blue}]{...} }
\lstset{literate=  {á}{{\'a}}1 {é}{{\'e}}1 {í}{{\'i}}1 {ó}{{\'o}}1 {ú}{{\'u}}1   {Á}{{\'A}}1 {É}{{\'E}}1 {Í}{{\'I}}1 {Ó}{{\'O}}1 {Ú}{{\'U}}1   {à}{{\`a}}1 {è}{{\`e}}1 {ì}{{\`i}}1 {ò}{{\`o}}1 {ù}{{\`u}}1   {À}{{\`A}}1 {È}{{\'E}}1 {Ì}{{\`I}}1 {Ò}{{\`O}}1 {Ù}{{\`U}}1   {ä}{{\"a}}1 {ë}{{\"e}}1 {ï}{{\"i}}1 {ö}{{\"o}}1 {ü}{{\"u}}1   {Ä}{{\"A}}1 {Ë}{{\"E}}1 {Ï}{{\"I}}1 {Ö}{{\"O}}1 {Ü}{{\"U}}1   {â}{{\^a}}1 {ê}{{\^e}}1 {î}{{\^i}}1 {ô}{{\^o}}1 {û}{{\^u}}1   {Â}{{\^A}}1 {Ê}{{\^E}}1 {Î}{{\^I}}1 {Ô}{{\^O}}1 {Û}{{\^U}}1   {œ}{{\oe}}1 {Œ}{{\OE}}1 {æ}{{\ae}}1 {Æ}{{\AE}}1 {ß}{{\ss}}1   {ű}{{\H{u}}}1 {Ű}{{\H{U}}}1 {ő}{{\H{o}}}1 {Ő}{{\H{O}}}1   {ç}{{\c c}}1 {Ç}{{\c C}}1 {ø}{{\o}}1 {å}{{\r a}}1 {Å}{{\r A}}1   {€}{{\euro}}1 {£}{{\pounds}}1 {«}{{\guillemotleft}}1   {»}{{\guillemotright}}1 {ñ}{{\~n}}1 {Ñ}{{\~N}}1 {¿}{{?`}}1 }
\usepackage{caption}
\usepackage{attachfile2}
\usepackage[margin=1.5cm,includeheadfoot,includehead,includefoot]{geometry}
\hypersetup{colorlinks,linkcolor=black}
\usepackage{fancyhdr}
\pagestyle{fancyplain}
\chead{}
\lhead{}
\rhead{}
\cfoot{}
\lfoot{\begin{footnotesize}alvaro.gonzalezsotillo@educa.madrid.org\end{footnotesize}}
\rfoot{\begin{footnotesize}\thepage / \pageref{LastPage}\end{footnotesize}}
\usepackage[skins]{tcolorbox}
\usepackage{multicol}
\usepackage{changepage} %ajdustwidth
\usepackage{fancybox}
\usepackage{attachfile2}
\lhead{Extraordinaria 2021 (es el \\lhead)}
\rhead{Administración y Gestión de Bases de Datos (es el \\rhead)}
\lhead{Seguridad y Alta Disponibilidad}
\rhead{Cortafuegos con iptables}
\usepackage{svg}
\usepackage{letltxmacro}
\LetLtxMacro{\originalincludegraphics}{\includegraphics}
\renewcommand{\includegraphics}[2][]{\IfFileExists{#2.pdf}{\originalincludegraphics[#1]{#2.pdf}}{\originalincludegraphics[#1]{#2}}}
\LetLtxMacro{\originalincludesvg}{\includesvg}
\renewcommand{\includesvg}[2][]{\IfFileExists{#2.pdf}{\originalincludegraphics[#1]{#2.pdf}}{\originalincludegraphics[#1]{#2.svg.pdf}}}
\usepackage{comment}
\excludecomment{NOTES}
\date{}
\title{Práctica 04: Cortafuegos con iptables en un router Docker}
\hypersetup{
 pdfauthor={},
 pdftitle={Práctica 04: Cortafuegos con iptables en un router Docker},
 pdfkeywords={},
 pdfsubject={},
 pdfcreator={},
 pdflang={Spanish}}
\begin{document}

\maketitle
\tableofcontents

\captionsetup{font=scriptsize}

\textattachfile{/agents/apuntes-clase_gh-pages-llm/seguridad-alta-disponibilidad/apuntes/4/sad-04-practica-cortafuegos-docker.org}{} \vspace{-1em}


\setlength{\parindent}{0em}
\setlength{\parskip}{1em}

\newtcolorbox{Aviso}[1][Aviso]{
  enhanced,
  colback=gray!5!white,
  colframe=gray!75!black,fonttitle=\bfseries,
  colbacktitle=gray!85!black,
  attach boxed title to top left={yshift=-2mm,xshift=2mm},
  title=#1
}

\newtcolorbox{cuadrito}[1][Ignorado]{
  %drop shadow southeast,
  enhanced jigsaw,
  colback=white,
}


\newcommand{\StudentData}{
  \begin{cuadrito}[1\textwidth]
    \vspace{0.3cm}
    \large{
      \textbf{Apellidos:} \hrulefill \\
      \textbf{Nombre:} \hrulefill \\
      \textbf{Fecha:} \hrulefill \hspace{1cm} \textbf{Usuario:} \hrulefill \\
    }
    \vspace{-0.2cm}
  \end{cuadrito}
}

\newcommand{\StudentDataDniFirma}{
  \begin{cuadrito}[1\textwidth]
    \vspace{0.3cm}
    \large{
      \textbf{Apellidos:} \hrulefill \\
      \textbf{Nombre:} \hrulefill  \\
      \textbf{Dni:} \hrulefill \hspace{8cm} \\
      \\
      \textbf{Firma:} \hrulefill \hspace{8cm} \\
    }
    \vspace{-0.2cm}
  \end{cuadrito}
}
\section{Objetivo de la práctica}
\label{sec:orgcfd6e54}

Al finalizar la práctica, el alumno será capaz de:

\begin{itemize}
\item Desplegar con \emph{Docker} una topología formada por una red externa, un router \emph{Linux} y una red interna aislada.
\item Comprobar que cada contenedor tiene su propia pila de red y sus propias tablas de \texttt{iptables}, independientes de las del host.
\item Publicar los servicios de la red interna hacia la red externa mediante \texttt{DNAT}.
\item Justificar y aplicar el \texttt{SNAT} necesario para que el tráfico publicado funcione en una red \texttt{internal} de \emph{Docker}.
\item Filtrar el tráfico con \texttt{iptables} usando políticas por defecto restrictivas y reglas de excepción.
\item Restringir el acceso a un servicio (\emph{MySQL}) a un único cliente autorizado.
\item Registrar en el log del kernel los intentos de acceso bloqueados y revisarlos.
\item Diagnosticar problemas de conectividad provocados por el cortafuegos.
\item Documentar las reglas y automatizar su aplicación.
\end{itemize}

Está disponible la \href{https://alvarogonzalezsotillo.github.io/apuntes-clase/seguridad-alta-disponibilidad/apuntes/4/sad-04-practica-cortafuegos-docker.pdf}{última versión de esta práctica en este enlace}.
\section{Descripción del escenario}
\label{sec:org50b1518}

Se parte del fichero \href{docker-compose.yml}{docker-compose.yml} incluido en esta carpeta. En él ya están definidos:

\begin{itemize}
\item La red \texttt{wan} (\texttt{172.30.0.0/24}), una red puente accesible desde el host, que representa la red real del host o la red externa del aula.
\item La red \texttt{lan} (\texttt{192.168.50.0/24}), declarada como \texttt{internal: true}, de modo que \emph{Docker} la aísla del exterior.
\item El contenedor \texttt{router} (\texttt{ubuntu:24.04}), con dos interfaces:
\begin{itemize}
\item \texttt{eth0} en \texttt{wan}, con IP \texttt{172.30.0.2}.
\item \texttt{eth1} en \texttt{lan}, con IP \texttt{192.168.50.2}.
\end{itemize}
\item El reenvío IP activado en el router.
\end{itemize}

El host tiene la IP \texttt{172.30.0.1} en la red \texttt{wan} y actúa como cliente externo. Los servicios que hay que añadir en la red interna son:

\begin{center}
\begin{tabular}{llrl}
Servicio & Contenedor & IP & Puerto\\
\hline
HTTP & \texttt{http} & 192.168.50.10 & 80/tcp\\
MySQL & \texttt{mysql} & 192.168.50.20 & 3306/tcp\\
\end{tabular}
\end{center}

Además se añadirá un cliente externo de administración:

\begin{center}
\begin{tabular}{llr}
Cliente & Contenedor & IP\\
\hline
Administración & \texttt{cliente-admin} & 172.30.0.100\\
\end{tabular}
\end{center}

El objetivo es que:

\begin{itemize}
\item Cualquier cliente de la red \texttt{wan} (por ejemplo, el propio host) pueda acceder al servidor HTTP mediante \texttt{http://172.30.0.2:8080}.
\item Solo \texttt{cliente-admin} (\texttt{172.30.0.100}) pueda acceder al servidor MySQL en \texttt{172.30.0.2:3306}.
\item Los intentos de acceso no autorizados a MySQL queden registrados en el log.
\item Ningún otro tráfico atraviese el router hacia la red interna.
\end{itemize}

\begin{center}
\includegraphics[width=.9\linewidth]{./media/practica-cortafuegos-docker.png}
\label{org3859277}
\end{center}

\textbf{Nota:} para que la práctica sea reproducible en cualquier equipo, la red \texttt{wan} se implementa como una red puente de \emph{Docker} accesible desde el host. Si se dispone de una interfaz física y se quiere simular la red real, se puede sustituir por una red \texttt{macvlan}; las reglas de \texttt{iptables} del router son las mismas.
\section{Requisitos}
\label{sec:orged0fd47}

\begin{itemize}
\item \emph{Docker} y \emph{Docker Compose} instalados.
\item Conexión a Internet para descargar las imágenes:
\begin{itemize}
\item \texttt{ubuntu:24.04}
\item \texttt{nginx:alpine} (servidor HTTP)
\item \texttt{mariadb:11} (servidor \emph{MySQL} compatible; la misma imagen incluye el cliente \texttt{mariadb})
\end{itemize}
\item Permisos para crear redes y contenedores. El router ya arranca con la capacidad \texttt{NET\_ADMIN} y el reenvío IP activado.
\item Un editor de texto y, en el host, un navegador o \texttt{curl}.
\end{itemize}
\section{Pasos a seguir}
\label{sec:orgfafef55}

\subsection{Completar el escenario}
\label{sec:orgff63caf}

Amplía el fichero \texttt{docker-compose.yml} para añadir los contenedores \texttt{http}, \texttt{mysql} y \texttt{cliente-admin}:

\begin{itemize}
\item \texttt{http}:
\begin{itemize}
\item Imagen \texttt{nginx:alpine}.
\item IP fija \texttt{192.168.50.10} en la red \texttt{lan}.
\item Una página \texttt{index.html} personalizada, montada en \texttt{/usr/share/nginx/html/index.html}.
\end{itemize}
\item \texttt{mysql}:
\begin{itemize}
\item Imagen \texttt{mariadb:11}.
\item IP fija \texttt{192.168.50.20} en la red \texttt{lan}.
\item Variables \texttt{MARIADB\_ROOT\_PASSWORD}, \texttt{MARIADB\_DATABASE}, \texttt{MARIADB\_USER} y \texttt{MARIADB\_PASSWORD}.
\end{itemize}
\item \texttt{cliente-admin}:
\begin{itemize}
\item Imagen \texttt{mariadb:11}, que ya incluye el cliente \texttt{mariadb}.
\item IP fija \texttt{172.30.0.100} en la red \texttt{wan}.
\item \texttt{entrypoint: ["sleep", "infinity"]}, para que no arranque ningún servidor y podamos usarlo como cliente.
\end{itemize}
\end{itemize}

Levanta el escenario y comprueba las redes y las IPs asignadas:

\begin{verbatim}
docker compose up -d
docker compose ps
docker network ls
docker network inspect 4_wan
docker network inspect 4_lan
\end{verbatim}

El prefijo de las redes (\texttt{4}) depende del nombre del directorio del proyecto.

Comprueba el punto de partida antes de filtrar nada:

\begin{verbatim}
# El router es accesible desde el host
ping -c 2 172.30.0.2

# Todavía no hay nada publicado en el puerto 8080 del router
curl -m 3 http://172.30.0.2:8080/
\end{verbatim}

La segunda orden falla con \emph{connection refused}: el router no está escuchando en el 8080, pero todavía \emph{contesta} con un rechazo. Más adelante, cuando el cortafuegos empiece a descartar paquetes, el fallo será por \emph{timeout}: conviene notar la diferencia.

Comprueba también que la red \texttt{lan} está realmente aislada para los clientes externos: \texttt{cliente-admin} no puede llegar a los servidores aunque conozca sus IPs.

\begin{verbatim}
docker exec cliente-admin mariadb --connect-timeout=3 \
  -h 192.168.50.10 -P 3306 -u root -pnoexiste -e "SELECT 1;"
\end{verbatim}

El cliente no consigue conectar: \emph{Docker} descarta en el host el tráfico que, procedente de otras redes, se dirige hacia un puente \texttt{internal}. Solo queda el camino del router, que es lo que vamos a controlar.
\subsection{Preparar el router}
\label{sec:org771a765}

Entra en el router e instala las herramientas que necesitarás para el cortafuegos y para diagnosticar:

\begin{verbatim}
docker exec -it router bash
apt-get update
apt-get install -y iptables iproute2 iputils-ping curl tcpdump conntrack
\end{verbatim}

Comprueba las interfaces, las rutas y el reenvío:

\begin{verbatim}
ip address
ip route
sysctl net.ipv4.ip_forward
\end{verbatim}

Observa las tablas de \texttt{iptables} del router:

\begin{verbatim}
iptables -L -n -v
iptables -t nat -L -n -v
\end{verbatim}

Están prácticamente vacías. Cada contenedor tiene su \emph{propia} pila de red y sus propias tablas, independientes de las del host. Compáralas con las que \emph{Docker} mantiene en el host (fuera del contenedor, con \texttt{sudo}):

\begin{verbatim}
sudo iptables -L -n -v
sudo iptables-save | grep -i docker
\end{verbatim}

En el host verás cadenas como \texttt{DOCKER}, \texttt{DOCKER-FORWARD}, \texttt{DOCKER-INTERNAL} y \texttt{DOCKER-USER} (en versiones anteriores, \texttt{DOCKER-ISOLATION-STAGE-1} y \texttt{DOCKER-ISOLATION-STAGE-2}). Son las que \emph{Docker} usa para gestionar sus puentes y aislar las redes internas, y conviene conocerlas porque explican parte del comportamiento del escenario.
\subsection{Políticas por defecto y tráfico básico}
\label{sec:orgafd5375}

Aplica políticas por defecto restrictivas y permite únicamente el tráfico imprescindible para el diagnóstico:

\begin{verbatim}
# Politicas por defecto
iptables -P INPUT DROP
iptables -P FORWARD DROP
iptables -P OUTPUT ACCEPT

# Loopback y conexiones ya establecidas
iptables -A INPUT -i lo -j ACCEPT
iptables -A INPUT -m conntrack --ctstate ESTABLISHED,RELATED -j ACCEPT
iptables -A FORWARD -m conntrack --ctstate ESTABLISHED,RELATED -j ACCEPT

# Diagnostico: permitir ping al router desde la red externa
iptables -A INPUT -i eth0 -p icmp --icmp-type echo-request -j ACCEPT
\end{verbatim}

Comprueba que el router sigue respondiendo a \texttt{ping} desde el host y que ahora el puerto 8080 ni siquiera rechaza la conexión:

\begin{verbatim}
ping -c 2 172.30.0.2
curl -m 3 http://172.30.0.2:8080/
\end{verbatim}

Con \texttt{INPUT} y \texttt{FORWARD} en \texttt{DROP}, ningún tráfico atraviesa todavía el router hacia la red interna.
\subsection{Publicar HTTP con DNAT}
\label{sec:org1232191}

Publica el servidor HTTP de la red interna en el puerto \texttt{8080} del router. Una primera aproximación razonable sería hacer solo \texttt{DNAT} y permitir el reenvío:

\begin{verbatim}
# Redirigir el trafico que llega al router por el puerto 8080
iptables -t nat -A PREROUTING -i eth0 -d 172.30.0.2 \
  -p tcp --dport 8080 -j DNAT --to-destination 192.168.50.10:80

# Permitir el trafico reenviado hacia el servidor HTTP
iptables -A FORWARD -i eth0 -o eth1 -p tcp \
  -d 192.168.50.10 --dport 80 \
  -m conntrack --ctstate NEW,ESTABLISHED,RELATED -j ACCEPT
\end{verbatim}

Comprueba desde el host:

\begin{verbatim}
curl -m 3 http://172.30.0.2:8080/
\end{verbatim}

\textbf{No funciona.} Diagnostica el problema con \texttt{tcpdump} en el router mientras repites el \texttt{curl} desde el host:

\begin{verbatim}
tcpdump -ni eth1 tcp port 80
\end{verbatim}

Ves los \texttt{SYN} saliendo por \texttt{eth1} hacia \texttt{192.168.50.10}, pero nunca vuelve ninguna respuesta: el paquete sale del router y no llega al servidor. La explicación está en el host: mira las reglas que \emph{Docker} instala para la red interna.

\begin{verbatim}
sudo iptables-save | grep -A5 DOCKER-INTERNAL
# o, en versiones anteriores de Docker:
sudo iptables-save | grep -A5 DOCKER-ISOLATION-STAGE-1
\end{verbatim}

Entre las reglas aparecen dos como estas (donde \texttt{br-XXXX} es el puente de la red \texttt{lan}):

\begin{verbatim}
-A DOCKER-INTERNAL ! -s 192.168.50.0/24 -o br-XXXX -j DROP
-A DOCKER-INTERNAL ! -d 192.168.50.0/24 -i br-XXXX -j DROP
\end{verbatim}

\emph{Docker} descarta en el host cualquier paquete que circule por el puente de una red \texttt{internal} con una IP de origen o de destino ajena a la subred. El paquete que hemos reenviado llega al puente con su origen original (\texttt{172.30.0.1}), así que el host lo descarta antes de que alcance al contenedor \texttt{http}.

La solución es que el router \emph{enmascare} ese tráfico con su propia IP de la red interna:

\begin{verbatim}
# SNAT del trafico que el router reenvia hacia la lan
iptables -t nat -A POSTROUTING -o eth1 -d 192.168.50.0/24 \
  -j SNAT --to-source 192.168.50.2
\end{verbatim}

Ahora el paquete entra en el puente con origen \texttt{192.168.50.2} (una IP de la subred), y la respuesta del servidor vuelve directamente al router, que restaura las IPs originales gracias a \texttt{conntrack}. Compruébalo:

\begin{verbatim}
curl -v http://172.30.0.2:8080/
\end{verbatim}

Debe aparecer la página personalizada. Comprueba también que el puerto 80 del router \emph{no} está publicado:

\begin{verbatim}
curl -m 3 http://172.30.0.2:80/
\end{verbatim}

Como efecto secundario del \texttt{SNAT}, el servidor HTTP ve las conexiones como si vinieran del router. Observa el log de accesos:

\begin{verbatim}
docker logs http
\end{verbatim}

Y revisa en el router las entradas de \texttt{conntrack} de la conexión que acabas de hacer:

\begin{verbatim}
conntrack -L | grep 192.168.50
\end{verbatim}

\textbf{Nota:} en un router físico, con los servidores usando al router como puerta de enlace, este \texttt{SNAT} no sería necesario. En este escenario hace falta precisamente por el aislamiento que \emph{Docker} aplica a las redes \texttt{internal}. Es el precio de que el host controle el acceso a la red interna y de que los servidores no necesiten ruta por defecto hacia el router: las respuestas vuelven a él por estar en la misma subred.
\subsection{Publicar MySQL con acceso restringido}
\label{sec:orge5cb692}

Publica MySQL en el puerto \texttt{3306} del router, pero permite el acceso únicamente al cliente de administración:

\begin{verbatim}
# Redirigir el trafico del puerto 3306 al servidor MySQL
iptables -t nat -A PREROUTING -i eth0 -d 172.30.0.2 \
  -p tcp --dport 3306 -j DNAT --to-destination 192.168.50.20:3306

# Permitir solo al cliente 172.30.0.100
iptables -A FORWARD -i eth0 -o eth1 -s 172.30.0.100 \
  -p tcp -d 192.168.50.20 --dport 3306 \
  -m conntrack --ctstate NEW,ESTABLISHED,RELATED -j ACCEPT

# Registrar y bloquear el resto de intentos
iptables -A FORWARD -i eth0 -o eth1 -p tcp \
  -d 192.168.50.20 --dport 3306 \
  -j LOG --log-prefix "MYSQL-DENEGADO: " --log-level 4
iptables -A FORWARD -i eth0 -o eth1 -p tcp \
  -d 192.168.50.20 --dport 3306 -j DROP
\end{verbatim}

Fíjate en el orden: la regla que acepta a \texttt{cliente-admin} debe estar \emph{antes} que las de \texttt{LOG} y \texttt{DROP}. El destino \texttt{LOG} no bloquea nada por sí mismo: solo registra el paquete y deja que siga evaluando reglas.

La regla de \texttt{SNAT} del paso anterior también se aplica a estas conexiones, así que el procedimiento es el mismo que con HTTP.

Comprueba el acceso desde \texttt{cliente-admin}:

\begin{verbatim}
docker exec cliente-admin mariadb --connect-timeout=5 \
  -h 172.30.0.2 -P 3306 -u alumno -palumno1234 \
  -e "SELECT VERSION();"
\end{verbatim}

Debe responder con la versión de \emph{MariaDB}. Comprueba desde el host que el acceso está bloqueado:

\begin{verbatim}
timeout 5 bash -c 'cat < /dev/null > /dev/tcp/172.30.0.2/3306'
echo $?
\end{verbatim}

El comando debe terminar con un código distinto de \texttt{0} y agotar el \emph{timeout} (un \texttt{DROP} no rechaza la conexión: la silencia). Revisa los intentos bloqueados en el log del kernel del host (el kernel es compartido entre el host y los contenedores, así que los mensajes de \texttt{iptables -j LOG} del router aparecen ahí):

\begin{verbatim}
sudo dmesg | grep MYSQL-DENEGADO
# o, si tu sistema usa journald:
sudo journalctl -k | grep MYSQL-DENEGADO
\end{verbatim}
\subsection{Guardar y automatizar las reglas}
\label{sec:org1546bc9}

Las reglas de \texttt{iptables} se pierden si el contenedor se reinicia. Reúnelas todas en un script \texttt{router/firewall.sh}, móntalo en el router con un volumen y ejecútalo tras levantar el escenario:

\begin{verbatim}
router:
  volumes:
    - ./router/firewall.sh:/firewall.sh:ro
\end{verbatim}

\begin{verbatim}
docker exec router bash /firewall.sh
\end{verbatim}

Como alternativa puedes volcar y restaurar las reglas con \texttt{iptables-save} e \texttt{iptables-restore} (fíjate en la opción \texttt{-i} para pasar el fichero por la entrada estándar):

\begin{verbatim}
docker exec router iptables-save > router/rules.v4
docker exec -i router iptables-restore < router/rules.v4
\end{verbatim}

Incluye el script en la entrega y explica qué harías para que las reglas se aplicaran automáticamente al recrear el escenario.
\subsection{Comprobaciones finales}
\label{sec:org413cfa9}

Rellena una tabla como esta con el resultado real de cada prueba:

\begin{center}
\begin{tabular}{lrll}
Origen & Destino & Puerto & Resultado esperado\\
\hline
host 172.30.0.1 & 172.30.0.2 & icmp & Permitido (ping)\\
host 172.30.0.1 & 172.30.0.2 & 8080/tcp & Permitido (HTTP)\\
host 172.30.0.1 & 172.30.0.2 & 80/tcp & Bloqueado\\
host 172.30.0.1 & 172.30.0.2 & 3306/tcp & Bloqueado (y en el log)\\
cliente-admin 172.30.0.100 & 172.30.0.2 & 3306/tcp & Permitido (MySQL)\\
cliente-admin 172.30.0.100 & 192.168.50.10 & cualquiera & Bloqueado (red aislada)\\
\end{tabular}
\end{center}

Incluye la salida de:

\begin{verbatim}
docker exec router iptables -L -n -v
docker exec router iptables -t nat -L -n -v
docker exec router conntrack -L
\end{verbatim}

y capturas de las pruebas permitidas y bloqueadas.
\section{Diagnóstico de problemas}
\label{sec:org648ccd6}

Responde brevemente a estas cuestiones:

\begin{enumerate}
\item Desde el host no se accede a \texttt{http://172.30.0.2:8080}. Describe cómo lo diagnosticarías con \texttt{tcpdump} en el router y qué revisarías en las tablas \texttt{nat} y \texttt{filter}.
\item \texttt{cliente-admin} no puede conectar con MySQL, pero el \texttt{ping} del host al router sí funciona. ¿Qué regla falta o está mal colocada?
\item Las conexiones al servidor HTTP llegan a \texttt{http} con dirección de origen \texttt{192.168.50.2}. ¿Por qué? ¿Qué implicación tiene esto para el servidor?
\item Aunque añadas una regla de \texttt{MASQUERADE} en el router, los contenedores de la \texttt{lan} no consiguen salir a Internet. Investiga con \texttt{sudo iptables-save} qué reglas instala \emph{Docker} en el host para las redes internas y explica el motivo.
\item Al reiniciar el contenedor \texttt{router} se pierden las reglas. ¿Cómo lo solucionas para que el escenario quede protegido nada más arrancar?
\item Comprueba si desde el host puedes acceder directamente a \texttt{http://192.168.50.10} sin pasar por el router, y compara con lo que le ocurre a \texttt{cliente-admin} al intentarlo. ¿Qué indica esto sobre dónde aplican las reglas de aislamiento de \emph{Docker} y por qué el host es un caso especial?
\end{enumerate}
\section{Qué se valorará}
\label{sec:org9d0769a}

Se entregará un documento (\texttt{PDF}, \texttt{ODT} o \texttt{DOCX}) con:

\begin{itemize}
\item \textbf{Topología:} diagrama con las redes, contenedores, IPs y flujo del tráfico.
\item \textbf{Despliegue:} \texttt{docker-compose.yml} completo y comandos utilizados.
\item \textbf{Cortafuegos:} script \texttt{firewall.sh} con las reglas comentadas y explicación de cada una, incluida la justificación del \texttt{SNAT}.
\item \textbf{Comprobaciones:} tabla de pruebas, salidas de \texttt{iptables} y \texttt{conntrack}, y capturas de los accesos permitidos y bloqueados.
\item \textbf{Logs:} captura de los intentos de acceso a MySQL bloqueados.
\item \textbf{Diagnóstico:} respuestas razonadas a los problemas planteados.
\item \textbf{Claridad y presentación:} pasos ordenados, explicaciones propias y capturas legibles.
\end{itemize}
\section{Instrucciones de entrega}
\label{sec:org9997052}

\begin{itemize}
\item \textbf{Modalidad:} individual. Solo se admitirá en pareja si así se indica expresamente en clase.
\item \textbf{Lugar de entrega:} tarea correspondiente a esta práctica en el aula virtual del módulo.
\item \textbf{Contenido:} documento de la práctica, \texttt{docker-compose.yml}, \texttt{firewall.sh} e \texttt{index.html}.
\item \textbf{Fecha:} la indicada en el aula virtual.
\end{itemize}
\end{document}
